\section{Recovering the curve subfields}
\label{sec:curve-recovery}

A \emph{curve subfield} of $K/k$ will mean a subfield $E\subset K$
containing $k$, relatively algebraically closed in $K$, and of
transcendence degree one over $k$. Such a field is finite over $k(t)$
for each $t\in E\setminus k$. We identify its completed multiplicative
group $\widehat E$ with its image in $\widehat K$ by
Lemma~\ref{prelim:completion}.

\begin{theorem}\label{recovery:curves}
The map $\Theta$ gives a bijection between the completed curve
subfields of $K/k$ and those of $L/l$. More precisely, for every curve
subfield $E\subset K$ there is a unique curve subfield $P\subset L$
such that
\[
 \Theta(\widehat E)=\widehat P.
\]
\end{theorem}

Theorem~\ref{pencil:uniform} bounds the degrees of the prime divisors
detected by a pencil. We first show that infinitely many occur. Their
incidence family gives containment in a completed curve field over a
finite extension of $L$. Intersections and norms then descend this
containment to $L$, and the inverse correspondence gives equality.

\subsection{The divisors detected by a pencil}

Fix $t\in K\setminus k$, let $E$ be the relative algebraic closure of
$k(t)$ in $K$, and put
\[
 h_a=\Theta[t-a],\qquad
 \mathcal S=\bigcup_{a\in k}\operatorname{Supp}_X(h_a).
\]
Here $X/l$ is the fixed normal projective model used in
Theorem~\ref{pencil:uniform}. That theorem bounds the degree of every
member of $\mathcal S$ by a constant depending on $t$ and $X$.

\begin{lemma}\label{recovery:residue-powers}
The set $\mathcal S$ is infinite. For every $z\in E^\times$, every
$j\geq1$, and every representative $g\in L^\times$ of
$\Theta[z]\bmod\ell^j$, all but finitely many $D\in\mathcal S$ satisfy
\begin{equation}\label{eq:recovery-residue-powers}
 \operatorname{ord}_D(g)=0,
 \qquad g|_D\in l(D)^{\times\ell^j}.
\end{equation}
\end{lemma}

\begin{proof}
The classes $[t-a]$ are independent in $k(t)^\times/\ell$, as evaluation
at their distinct zeros shows. The kernel of
$k(t)^\times/\ell\to K^\times/\ell$ is finite dimensional: adjoining
roots for any finite independent subset of that kernel gives an
elementary abelian $\ell$-extension contained in the finite extension
$E/k(t)$. Thus their images span an infinite-dimensional subspace of
$K^\times/\ell$, and the same is true of the $h_a$ in
$\widehat L/\ell=L^\times/\ell$.

The kernel of divisor evaluation
\[
 L^\times/\ell\longrightarrow
 \bigoplus_{D\in X^{(1)}}\mathbb F_\ell
\]
is finite dimensional. Indeed, take a smooth dense open $U\subset X$.
A function whose divisor on $U$ is divisible by $\ell$ defines a
$\mu_\ell$-torsor on $U$: its divisor divided by $\ell$ is Cartier.
The kernel therefore embeds in
$H^1_{\mathrm{\acute et}}(U,\mu_\ell)$, which is finite dimensional
by Deligne's finiteness theorem
\cite[Th\'eor\`emes de finitude, Theorem 1.1]{DeligneSGA4half}.
If $\mathcal S$ were finite, both the image and the kernel
of divisor evaluation on the pencil span would be finite dimensional,
a contradiction.

For $D\in\mathcal S$, let $w_D$ be its divisorial valuation and let $v$
be the quasi-prime divisor of $K/k$ supplied by
Proposition~\ref{local:transport}. Its inertia detects some $[t-a]$.
Consequently $vE/vk$ is nonzero inside the rank-one group $vK/vk$.
The valuation transcendence-degree inequality applied to $E/k$ gives
$\operatorname{trdeg}(Ev/kv)=0$, so $Ev=kv$ because $kv$ is
algebraically closed.

If $\operatorname{ord}_D(\Theta[z])=0$, inertia evaluation implies
$v(z)\in vk$. Multiplying $z$ first by a constant to give it value zero,
and then by a constant with the reciprocal residue, makes it a
principal unit. Thus $D_v$ kills $[z]$, and $D_{w_D}$ kills $\Theta[z]$.
By Corollary~\ref{pencil:finite}, the excluded $D$ form a finite set.
Discard also the finitely many prime divisors in the ordinary divisor
of $g$. For every remaining $D$, $g$ is a unit at $D$, and its residue
has zero evaluation modulo $\ell^j$ against
$D_{w_D}/I_{w_D}=W_{l(D)}$. Finite-level Kummer duality gives
\eqref{eq:recovery-residue-powers}.
\end{proof}

\subsection{One incidence curve for every Kummer test}

The following geometric lemma permits an arbitrary collection of tests.
In particular, neither the constants nor the collection must be
countable or uncountable.

\begin{lemma}\label{recovery:incidence}
Let $l$ be algebraically closed, let $X/l$ be a normal integral projective
variety of dimension $d\geq2$, and let $\mathcal S$ be an infinite set
of prime divisors of bounded degree in a fixed projective embedding.
Let $\mathcal A$ be any collection of pairs $(g,q)$, where
$g\in l(X)^\times$ and $q\geq1$ is prime to $\operatorname{char}l$.
Suppose that, for each pair separately, all but finitely many
$D\in\mathcal S$ satisfy
\[
 \operatorname{ord}_D(g)=0,
 \qquad g|_D\in l(D)^{\times q}.
\]
There exist a finite field extension $M/l(X)$ and a curve field
$P_0/l$ contained in $M$ such that $M/P_0$ is regular and
\[
 g\in(M\overline{P_0})^{\times q}
 \quad\text{for every }(g,q)\in\mathcal A.
\]
The same fields $M$ and $P_0$ work for the entire collection.
\end{lemma}

\begin{proof}
\emph{Place the divisors in a bounded family.}
Integral subvarieties of fixed dimension and bounded degree in
$\mathbb P^N_l$ have only finitely many Hilbert polynomials
\cite[Section~2 and Lemma~2.4]{GrothHilbert1961}; see also
\cite[Proposition~5.3]{Tanaka2023}.

The Hilbert points of $\mathcal S$ therefore lie in a finite union of
projective Hilbert schemes of $X$, with their universal flat families
\cite[Theorem~3.2 and the Hilbert specialization after Proposition~3.8]{GrothHilbert1961}.
Their closure has a positive-dimensional
irreducible component on which they are dense. Give it its reduced
structure and shrink to an integral variety $T$ such that the universal
family
\[
 Z\subset T\times X\longrightarrow T
\]
has geometrically integral divisor fibers. This is possible by
constructibility of geometric integrality
\cite[Tags 0579, 055B]{Stacks}; the family is flat by construction.

\emph{Fix one incidence curve.}
Choose an integral locally closed curve $T'\subset T$ now, before
making any of the tests in $\mathcal A$, and put $Z'=Z\times_T T'$.
Both $Z$ and $Z'$ are integral. Indeed, flatness over an integral base
excludes vertical components and embeds affine coordinate rings in
their generic-fiber localizations, while the generic fibers are
geometrically integral. Both incidence maps to $X$ are dominant.
Otherwise a proper closed subset of $X$ would contain infinitely many
distinct prime-divisor fibers; each would have to be one of its
finitely many components of dimension $d-1$. Distinct parameter points
give distinct divisors because they are distinct Hilbert points.

Set
\[
 M=l(Z'),\qquad P_0=l(T').
\]
Dominance embeds $l(X)$ into $M$. Since $\dim Z'=d$, this is a finite
extension. Geometric integrality of the generic fiber of $Z'\to T'$
means that $M/P_0$ is regular.

Figure~\ref{fig:incidence} records the two projections that determine
these fields.
\begin{figure}[htbp]
\centering
\begin{tikzcd}[column sep=large,row sep=large]
 Z' \arrow[r] \arrow[d] & Z \arrow[r,"\operatorname{pr}_X"] \arrow[d]
   & X \\
 T' \arrow[r,hook] & T &
\end{tikzcd}
\caption{The left square is Cartesian. The composite $Z'\to X$ is
dominant and generically finite, giving $l(X)\subset M=l(Z')$;
the vertical map gives $P_0=l(T')\subset M$. The curve $T'$ is fixed
before the tests $(g,q)$. Auxiliary normal covers used to split a test
may depend on that test, but $M$ and $P_0$ do not.}
\label{fig:incidence}
\end{figure}

\emph{Extend each splitting to the fixed curve.}
It remains to prove all the power assertions without changing $T'$.
Fix one pair $(g,q)$. Let $U\subset X$ be the unit locus of $g$, and let
\[
 Y=Z^{\mathrm{sm}/T}\cap\operatorname{pr}_X^{-1}(U).
\]
Then $Y\to T$ is smooth. Its fibers over the generic points of $T$
and of $T'$ are nonempty and geometrically integral: geometric
integrality supplies a dense smooth locus, and dominance of the two
incidence maps supplies the unit condition. Over $Y$ consider the
degree-$q$ finite \'etale cover
\[
 C_g=\operatorname{Spec}_Y\mathcal O_Y[u]/(u^q-g).
\]
On every tested fiber satisfying the hypothesis and meeting $Y$, a
rational $q$-th root extends to a regular unit by normality of the
smooth fiber. Thus the cover splits into $q$ copies on those fibers.
The number of geometric irreducible components is constant on a dense
open of $T$ \cite[Tag 055A]{Stacks}. That number must be $q$, by the
dense set of tested fibers. A degree-$q$ finite \'etale cover of a
normal integral scheme with $q$ components has degree one on each component.
It follows that the cover is split on the geometric generic fiber
over $T$.

This last dense open need not meet $T'$. To pass to the fixed curve,
descend the generic splitting to a finite extension of $l(T)$ and
normalize $T$ in it, obtaining a finite surjective morphism
$\widetilde T\to T$. The scheme
$\widetilde Y=Y\times_T\widetilde T$ is normal because it is smooth
over a normal base \cite[Tag 034F]{Stacks}. It is integral because its
generic fiber is geometrically integral and it is flat over the
integral base. The pulled-back cover is generically split and finite
\'etale over this normal integral scheme, hence split everywhere:
each component is finite birational over the normal target and is an
isomorphism \cite[Lemma 58.11.2, Tag 0BQL]{Stacks}.
Surjectivity of $\widetilde T\to T$ provides a lift of a geometric
point above the generic point of $T'$. Restricting the split cover to
that point proves that $g$ has a $q$-th root in
$M\overline{P_0}$. The auxiliary finite extension may depend on
$(g,q)$; the already chosen fields $M$ and $P_0$ do not.
\end{proof}

\subsection{Kummer descent and intersections}

We next turn geometric splitting into membership in a completed curve
field. Two elementary facts keep track of the finite extension
introduced by the incidence construction.

\begin{lemma}\label{recovery:kummer-descent}
Let $M/P$ be a regular field extension, let $q$ be prime to the
characteristic, and suppose $\mu_q\subset P$. If
$g\in M^\times$ becomes a $q$-th power in $M\overline P$, then its
class in $M^\times/M^{\times q}$ comes from a unique class in
$P^\times/P^{\times q}$.
\end{lemma}

\begin{proof}
A $q$-th root $r$ of $g$ is separable over $M$, so it already belongs to
$MP^{\mathrm{sep}}$: the remaining constant extension is purely
inseparable. Choose a finite Galois extension $P'/P$ with $r\in MP'$.
Regularity identifies
$\operatorname{Gal}(MP'/M)$ with $\operatorname{Gal}(P'/P)$.
The ratios $\sigma(r)/r\in\mu_q$ form a character of this group.
By Hilbert's Theorem~90, there is $b\in P'^\times$ with these same
ratios. Then $r/b\in M$ and $b^q\in P$, giving
$g=b^q(r/b)^q$. Finally, if an element of $P$ has a $q$-th root in
$M$, that root is algebraic over $P$ and therefore belongs to $P$.
This proves uniqueness.
\end{proof}

\begin{lemma}\label{recovery:intersection}
Let $N/l$ be a function field over algebraically closed constants of
characteristic different from $\ell$. If $P$ and $Q$ are distinct
curve subfields of $N/l$, then
$\widehat P\cap\widehat Q\subset\widehat N$ has finite
$\Zell$-rank.
\end{lemma}

\begin{proof}
The compositum $PQ$ has transcendence degree two over $l$; otherwise
relative algebraic closedness in $N$ would force $P=Q$. Let $C_P$ and
$C_Q$ be the smooth projective curves of these fields. The product is
integral, and its function field identifies with $PQ$: the rational
map defined by the two subfields has two-dimensional, hence dense,
image in $C_P\times C_Q$.

For a point $x\in C_P(l)$, take the valuation of the divisor
$\{x\}\times C_Q$. Extend it by a Gauss valuation over a transcendence
basis for $N/PQ$, then prolong it across the remaining finite
extension. Its normalized order $w$ on $N$ is zero on $Q^\times$ and
restricts to $e\operatorname{ord}_x$ on $P^\times$, for some integer
$e>0$. This construction allows inseparable extensions and all
ramification indices.

An element of $\widehat P\cap\widehat Q$ has $w$-order zero by its
$Q$-representation. Its $x$-order is therefore zero, since
multiplication by $e$ is injective on $\Zell$. This holds for every
$x$. By Lemma~\ref{prelim:curve-kernel}, the intersection embeds in
the finite-rank Tate module of $\operatorname{Pic}^0(C_P)$.
\end{proof}

\begin{proof}[Proof of Theorem~\ref{recovery:curves}]
Let $E\subset K$ be a curve subfield, and choose $t\in E\setminus k$.
Relative algebraic closedness makes $E$
the relative algebraic closure of $k(t)$ in $K$. Apply
Lemmas~\ref{recovery:residue-powers} and \ref{recovery:incidence} to
one representative of each $\Theta[z]\bmod\ell^j$, for
$z\in E^\times$ and $j\geq1$. We obtain fields $L\subset M$ and
$P_0\subset M$ as in the incidence lemma.

Put $H=\Theta(\widehat E)$. We claim that its image in $\widehat M$
lies in $\widehat P_0$. At level $\ell^j$, every element of
$\widehat E$ is represented by an actual $z\in E^\times$.
Lemma~\ref{recovery:kummer-descent} puts its image in
$P_0^\times/P_0^{\times\ell^j}$. These classes are unique because
$P_0$ is relatively algebraically closed in $M$, and their uniqueness
makes them compatible as $j$ varies. Taking inverse limits proves
the claim. Completion maps across $M/L$ are injective by
Lemma~\ref{prelim:completion}.

Choose a finite normal extension $N/L$ containing $M$, allowing an
inseparable part, and let $P_N$ be the relative algebraic closure of
$P_0$ in $N$. It is a curve subfield of $N/l$. The image of $H$ lies
in $\widehat P_N$. Every $\sigma\in\operatorname{Aut}_L(N)$ fixes
$H$, so $H$ also lies in $\widehat{\sigma(P_N)}$.
It has infinite rank: $\widehat E/\ell$ has infinite dimension, as
the pencil argument in Lemma~\ref{recovery:residue-powers} shows.
Indeed, lifts of any finite mod-$\ell$ independent set are
$\Zell$-independent: a nonzero relation, divided by the smallest
$\ell$-power in its coefficients, would reduce to a nonzero
mod-$\ell$ relation. Torsion freeness justifies this division.
Lemma~\ref{recovery:intersection} therefore forces
$\sigma(P_N)=P_N$ for every $\sigma$.

Set $P_L=P_N\cap L$. This field has transcendence degree one over
$l$. Indeed, for $G=\operatorname{Aut}_L(N)$, the field $P_N^G$ has
transcendence degree one and lies in $N^G$. In characteristic zero
$N^G=L$. In characteristic $p>0$, $N^G/L$ is finite purely inseparable,
so a common $p^a$-th power sends $P_N^G$ into $P_L$. Moreover $P_L$
is relatively algebraically closed in $L$: an element of $L$
algebraic over $P_L$ is algebraic over $P_N$, hence lies in $P_N$.
Taking the relative algebraic closure of any nonconstant rational
subfield shows that $P_L/l$ is finitely generated.

The norm $N_{N/L}$ sends $P_N^\times$ into $P_L^\times$. In fact,
if $s$ is the inseparable degree of $N/L$, then
\[
 N_{N/L}(x)=\left(\prod_{\sigma\in G}\sigma(x)\right)^s;
\]
the right side lies in $P_N$ because $P_N$ is $G$-stable, and the norm lies in
$L$. On completions, norm composed with inclusion is multiplication
by $[N:L]$. We conclude that
$[N:L]H\subset\widehat P_L$. The latter is saturated in
$\widehat L$ by Lemma~\ref{prelim:completion}, so
$H\subset\widehat P_L$.
Here the prime-to-$\ell$ factor of $[N:L]$ is a unit of $\Zell$,
and saturation removes its remaining $\ell$-power factor.

Apply the same argument to $\Theta^{-1}$. It gives a curve subfield
$E'\subset K$ with
\[
 \widehat E\subset\Theta^{-1}(\widehat P_L)
 \subset\widehat E'.
\]
The infinite-rank intersection forces $E'=E$ by
Lemma~\ref{recovery:intersection}. Both inclusions are therefore
equalities. That lemma also proves uniqueness. Reversing $K$ and $L$
proves the asserted bijection.
\end{proof}

\subsection{Divisorial valuations and rational quotients}

We can now distinguish valuations trivial on the constants. Throughout
this subsection, a \emph{divisorial valuation} means such a
quasi-prime divisor, with its discrete value group normalized to
$\mathbb Z$.

The following criterion adapts Topaz's Lemma~A.7
\cite[Lemma~A.7]{TopazTorelli2025} to completed multiplicative groups.
The proof below verifies the needed hypotheses in this setting.

\begin{proposition}\label{recovery:divisorial}
A quasi-prime divisor $v$ of $K/k$ is divisorial if and only if
\begin{equation}\label{eq:recovery-divisorial}
 D_v^\perp\cap\widehat E=0
\end{equation}
for some curve subfield $E$ of $K/k$. Consequently $\varphi$
preserves divisorial inertia and decomposition groups in both
directions.
\end{proposition}

\begin{proof}
Suppose $v$ is trivial on $k$. Its residue field has transcendence
degree $d-1\geq1$. Choose $t$ with transcendental residue, and let
$E$ be the relative algebraic closure of $k(t)$ in $K$. The valuation
is trivial on $k(t)$ and hence on its algebraic extension $E$.
Thus $E$ embeds into $Kv$. Its relative algebraic closure in $Kv$
is finite over $E$, so the induced completion map is injective:
factor it through that finite extension and then a relatively
algebraically closed inclusion. Residual evaluation by
$D_v/I_v=W_{Kv}$ proves \eqref{eq:recovery-divisorial}.

Conversely, suppose $v|k$ is nontrivial, and let $E$ be any curve
subfield. Choose $x\in E\setminus k$, replacing it by $x^{-1}$ if
necessary so that $v(x)\geq0$, and choose $c\in k$ with $v(c)>0$.
Then $1+cx$ is a nonconstant principal unit. Its class in
$\widehat E$ is nonzero and is killed by $D_v$, so
\eqref{eq:recovery-divisorial} fails. The criterion is invariant
under $\Theta$ by Theorem~\ref{recovery:curves} and the local
correspondence.
\end{proof}

For a curve subfield $E\subset K$, restriction gives a surjection
\[
 \rho_E:W_K\longrightarrow W_E.
\]
This is the continuous surjection of Lemma~\ref{prelim:completion}.

\begin{proposition}\label{recovery:rational}
The point inertia lines of every curve subfield $E$ are precisely
the saturations of the nonzero images $\rho_E(I_v)$ for divisorial
valuations $v$ of $K/k$. The correspondence of
Theorem~\ref{recovery:curves} preserves these lines and the genus.
In particular, it matches exactly the relatively algebraically closed
rational subfields.

If $\Theta(\widehat E)=\widehat P$, the induced continuous
isomorphism $\varphi_E:W_P\to W_E$ makes the diagram
\begin{equation}\label{eq:recovery-quotient-diagram}
\begin{tikzcd}[column sep=large]
 W_L \arrow[r,"\varphi"] \arrow[d,two heads,"\rho_P"']
   & W_K \arrow[d,two heads,"\rho_E"] \\
 W_P \arrow[r,"\varphi_E"'] & W_E
\end{tikzcd}
\end{equation}
commute and matches the point inertia lines on its targets.
For rational $E$ and $P$, these are the geometric rational quotient
diagrams, with their point decomposition data.
\end{proposition}

\begin{proof}
A divisorial valuation restricted nontrivially to $E$ is an integer
multiple of a point valuation on its smooth projective curve: its
value group is a nonzero subgroup of $\mathbb Z$, and it is trivial
on $k$. Conversely, extend any point valuation of $E$ by a Gauss
valuation over a transcendence basis for $K/E$, then prolong it over
the remaining finite extension. The resulting valuation is
divisorial on $K/k$. Its restriction is $e\operatorname{ord}_x$ for
some $e>0$. Saturation of the nonzero image recovers
$\Zell\operatorname{ord}_x$, including when $\ell$ divides $e$.

The dual of $\Theta|_{\widehat E}$ defines $\varphi_E$.
Evaluation against elements of $\widehat E$ proves the commutativity
of \eqref{eq:recovery-quotient-diagram}. Proposition~\ref{recovery:divisorial}
and the preceding description of point lines show that $\varphi_E$
matches them. Their generators may differ by units, which do not
change the common kernel of point-order evaluations modulo $\ell$.
Its dimension is twice the genus by Lemma~\ref{prelim:curve-kernel}.
Hence genus is preserved. A smooth projective curve over an
algebraically closed field has genus zero exactly when its function
field is rational.

If $k(t)$ is relatively algebraically closed in $K$, the extension
$K/k(t)$ is regular. Only separability needs explanation. In
characteristic $p>0$, relative algebraic closedness implies
$t\notin K^p$. Since $k$ is perfect, $t$ can be included in a
separating transcendence basis for $K/k$, proving that $K/k(t)$ is
separably generated. Characteristic zero is immediate. Thus these
are exactly the rational subfields defined by general elements.
Finally, the residue field at a point of a rational curve is the
algebraically closed constant field, whose character group is zero.
Point decomposition therefore equals point inertia, and the diagram
also carries the required point decomposition data.
\end{proof}

The curve correspondence has consequently recovered both the true
divisorial valuations and all geometric rational quotients from the
original datum. The final reconstruction theorem can now be applied
with these intrinsically determined objects.
